% AI reflection (half a page, ungraded; hand-in with the paper and slides, Fri 27 Nov 2026, DTU Learn).
% The course question: did the "From an AI-drafted report" boxes change how you read your draft, and where did you
% use them (or not)? Final text by Martin, Anas and Diego (27 Nov 2026). Compile with pdfLaTeX; fits one A4 page.
\documentclass[11pt,a4paper]{article}
\usepackage[utf8]{inputenc}
\usepackage[T1]{fontenc}
\usepackage{lmodern}
\usepackage[margin=2.5cm]{geometry}
\usepackage{microtype}
\setlength{\parindent}{0pt}
\setlength{\parskip}{0.8em}
\pagestyle{empty}

\begin{document}

{\large\bfseries AI reflection}\\[0.2em]
{\small Anas Mir, Martin M{\o}llenhus and Diego A. Mijares \quad|\quad 02805 Social Graphs and Interactions}

\bigskip

The boxes changed how we read our draft by making us look at it critically. They show how confidently an LLM can
state a wrong answer. It has no inner critic, and when its information is incomplete it infers from what it knows
instead of asking questions. So wherever AI helped with the report, we made sure it had the full context, asked us
before doing inference and by us critically reviewing the draft.

The boxes also show what drift feels like. An AI and the person using it can slowly stop sharing an understanding of
what is being said and why, and the AI keeps going anyway. It skips ahead, and the sentence comes out ``fluent'' and
confident, like the overclaiming paragraphs in the boxes. A paper like ours (or any paper for that matter) needs
checks behind every claim, so drift is the easiest way to say something we have not earned. We tried to build the
guardrails in from the start instead of catching problems at the end. Claude worked from written rules, a log where
every modelling choice needed our reasons before it counted, and a record of every place a model touched a result.
Each of us prompted our Claude Code to build the context of what we did into one shared repository and to scrutinize
it and us, and we verified it by checking each other's work there. We designed how the agents go through our workflow
so that Claude acts like a critical scientist, aimed at sharpening our thinking. When it ran ahead, we pulled it back
to the question and read its work line by line.

The same setup made hand-labelling cheaper. Every time we changed the name matcher, checking it by hand took too much
effort. So we had several Claude sessions with no context read the text blind, compared their answers with the
matcher's, and read only the cases where they disagreed. That let us run many more checks and left time for the ideas
that matter.

The model writes fast. Staying in step with it was our job.

\end{document}
